\documentclass{scrartcl}
\usepackage{amsmath}
\usepackage{amssymb}

\title{maj2random}
\author{\small Michael Clark \\ \small Independent Researcher \\ \small \texttt{<michaeljclark@mac.com>}}
\date{}

\begin{document}
\maketitle

\subsection*{introduction}

\noindent
\emph{maj2random} is a low-overhead floating point random number generator
built upon a reduced-round variant of the SHA-2 compression function. it
retains the high-quality mixing behavior of SHA-2 to permute a \emph{GLSL
vec2} input, producing two pseudorandom values in the interval $[0,1]$.
the name \emph{maj2random} derives from the use of the \texttt{maj} mixing
function and pays homage to \texttt{arc4random}; the “2” refers to the number
of rounds.
\\
\par\noindent
cryptographically secure hash functions are useful building blocks for
pseudorandom number generators because they are statistically diffuse,
whereas general-purpose hash functions are optimized for collision
resistance but typically exhibit biased statistics.
\\
\par\noindent
the goal is a pseudorandom number generator that trades cryptographic
strength for speed, making it suitable for GPU shaders and simulations where
high-quality statistically diffuse randomness is needed but where the full
SHA-256 hash function is too costly. the algorithm produces values
closely approximating a uniform distribution on $[0,1]$, using only 2 rounds
at $\approx 3\%$ of the overhead of the full SHA-256 hash function.

\subsection*{functions}

\noindent
the SHA-256 algorithm uses six functions, with each function
operating on 32-bit words, denoted \(x\), \(y\), and \(z\).
\textit{maj2random} adopts these functions unchanged and adds \emph{unorm}
to compute a normalized floating-point sample from an integer:

\[
\begin{array}{rclclcl}
\mathrm{ch}(x,y,z)  & = & \multicolumn{5}{l}{(x \land y) \oplus (x \land z)} \\
\mathrm{maj}(x,y,z) & = & \multicolumn{5}{l}{(x \land y) \oplus (x \land z) \oplus (y \land z)} \\[0.5em]
\Sigma_0(x) & = & \mathrm{ror}(x,2)  & \oplus & \mathrm{ror}(x,13) & \oplus & \mathrm{ror}(x,22) \\
\Sigma_1(x) & = & \mathrm{ror}(x,6)  & \oplus & \mathrm{ror}(x,11) & \oplus & \mathrm{ror}(x,25) \\
\sigma_0(x) & = & \mathrm{ror}(x,7)  & \oplus & \mathrm{ror}(x,18) & \oplus & \mathrm{shr}(x,3) \\
\sigma_1(x) & = & \mathrm{ror}(x,17) & \oplus & \mathrm{ror}(x,19) & \oplus & \mathrm{shr}(x,10) \\[0.5em]
\mathrm{unorm}(n) & = & \frac {n \; \land \; (2^{23}-1)}{2^{23}-1}
\end{array}
\]

\newpage
\subsection*{constants}

\noindent
the SHA-256 algorithm uses sixty-four constant 32-bit words $K_{0...63}$
containing the first 32-bits of the fractional part of the cube roots of the
first sixty-four prime numbers. \textit{maj2random} uses only the first two.
the first eight are included for testing purposes.

\[
K_i = \begin{bmatrix}
    \mathtt{0x428a2f98}, & \mathtt{0x71374491}, & \mathtt{0xb5c0fbcf}, & \mathtt{0xe9b5dba5}, \\
    \mathtt{0x3956c25b}, & \mathtt{0x59f111f1}, & \mathtt{0x923f82a4}, & \mathtt{0xab1c5ed5}
\end{bmatrix}, \quad 0 \le i < 8
\]

\vspace{1em}

\subsection*{entropy}

\noindent
\textit{maj2random} requires 48-bits of entropy extracted from two single
precision \emph{IEEE754.1} floating-point numbers such as a \emph{GLSL vec2},
with implicit leading one in the \emph{mantissa}:

\[
\begin{array}{rcl}
uv & = &
\begin{bmatrix}
  (-1)^{\mathit{sign}_x} \cdot \mathrm{mantissa}_x \cdot 2^{\mathit{exponent}_x} \\
  (-1)^{\mathit{sign}_y} \cdot \mathrm{mantissa}_y \cdot 2^{\mathit{exponent}_y}
\end{bmatrix}
\end{array}
\]

\vspace{1em}

\noindent
\textit{maj2random} extracts 48 bits of entropy from the \emph{sign} and
\emph{mantissa} fields of $uv$. these bits are combined using funnel-shift
operations and XOR mixing to form the first two words of the message vector
$M_i$. the exponent is ignored as inputs are expected to be normalized $uv$
coordinates (e.g.\ texture coordinates)\footnote{exponent entropy could be
included, but earlier \emph{GLSL} versions lack reliable \texttt{frexp} or
\texttt{ilogb}, and \texttt{log2} is insufficiently precise.}.

\[
\begin{array}{rcl}
x & = & \mathrm{uint32}\Big( \mathrm{abs}(\mathit{uv.x}) \cdot 2^{23} \lor \;
        \mathrm{sign}(\mathit{uv.x}) \cdot 2^{23} \Big) \\
y & = & \mathrm{uint32}\Big( \mathrm{abs}(\mathit{uv.y}) \cdot 2^{23} \lor \;
        \mathrm{sign}(\mathit{uv.y}) \cdot 2^{23} \Big) \\[1em]
M_i & = &
\begin{bmatrix}
x \; \oplus \; (\mathrm{shl}(x, 8) \; \lor \; \mathrm{shr}(y, 16)), &
y \; \oplus \; (\mathrm{shl}(x, 16) \; \lor \; \mathrm{shr}(y, 8)))
\end{bmatrix}
\end{array}
\]

\vspace{1em}

\subsection*{key schedule}

\noindent
the key schedule expands the initial entropy into a full message block
$W_0\dots W_n$ suitable for hashing. in \textit{maj2random}, because
only two message words are available, the 64 bits of state from
$M_0\dots M_1$ are replicated over $W_i$ modulo 2:

\[
W_i =
\begin{array}{ll}
M_{i \bmod 2}.
\end{array}
\]

\vspace{1em}

\noindent
\emph{note:} the regular SHA-2 key schedule is omitted because the input
entropy is premixed.

\newpage
\subsection*{algorithm}

\noindent
the standard SHA-256 compression function mixes the message block
$W_0\dots W_i$ into the internal state $H_0\dots H_7$ to produce
pseudo-random outputs. in \textit{maj2random}, only 2 rounds are applied,
and the initial state $H_0\dots H_7$ is set to \textit{zero} beforehand.

\[
\text{for } i = 0,1:
\quad
\begin{aligned}
T_0  &= W_i + H_7 + \sigma_1(H_4) + \mathrm{ch}(H_4, H_5, H_6) + K_i \\
T_1  &= \mathrm{maj}(H_0, H_1, H_2) + \sigma_0(H_0) \\
H_7 &= H_6, \; H_6 = H_5, \; H_5 = H_4 \\
H_4 &= H_3 + T_0 \\
H_3 &= H_2, \; H_2 = H_1, \; H_1 = H_0 \\
H_0 &= T_0 + T_1
\end{aligned}
\]

\vspace{1em}

\subsection*{sample}

\noindent
the \emph{unorm} function is applied to the XOR of the hash state
$H_0\dots H_7$ produced by the compression function to compute $S_i$, a
floating-point sample in $[0,1]^2$:

\[
S_i =
\begin{bmatrix}
S_x \\[1mm] S_y
\end{bmatrix}
=
\begin{bmatrix}
\mathrm{unorm}(H_0 \oplus H_1 \oplus H_2 \oplus H_3) \\[1mm]
\mathrm{unorm}(H_4 \oplus H_5 \oplus H_6 \oplus H_7)
\end{bmatrix}
\]

\vspace{1em}

\subsection*{statistics}

\noindent
mean, variance, standard deviation, and chi-squared ratio are measured for
\emph{maj2random} using $(x,y)$ coordinates composed of a combination of
random values and zeroes for $\mathtt{NROUNDS} \in \{2, 4, 8\}$. the standard
deviations range from $[0.28757,\; 0.28865]$ across all combinations,
compared with the theoretical value for a uniform distribution on \([0,1]\),
\[
  \sigma = \frac{1}{\sqrt{12}} \approx 0.288675.
\]

\vspace{1em}

\noindent
a normalized $\chi^2$ (chi-squared) statistic, the \emph{chi-squared ratio}
$\chi^2_\mathrm{ratio}$, was computed as:
\[
  \chi^2_\mathrm{ratio} = \frac{(n-1)s^2}{n\,\sigma^2_\mathrm{theory}} =
    \frac{s^2}{\sigma^2_\mathrm{theory}} \, \frac{n-1}{n},
\]

\vspace{1em}

\noindent
which expresses the measured variance relative to the theoretical variance.
$\chi^2_\mathrm{ratio}$ values fall within  $[0.99252,\,0.99972]$,
$[0.99237,\,0.99985]$, and $[0.99244,\,0.99950]$  for $\mathtt{NROUNDS}
\in \{2, 4, 8\}$, respectively. this indicates that the  observed variances
are extremely close to the theoretical variance, with minimal difference
between $\mathtt{NROUNDS}=2$ and $\mathtt{NROUNDS}=8$. values near $1.0$
correspond to perfect agreement with the theoretical uniform distribution,
and there is no evidence of statistically significant deviation.
\newpage

\begin{table}[h]
\centering
\small
\begin{tabular}{llrrrrr}
\hline
\textbf{x-coord} & \textbf{y-coord} & \textbf{count} & \textbf{mean} &
\textbf{variance} & \textbf{std-dev} & \textbf{$\chi^2$ ratio} \\
\hline
          0 & (0 - 1M)/8M  & 1,000,000 & 0.50015 & 0.08320 & 0.28844 & 0.99835 \\
          0 & (0 - 8M)/8M  & 8,000,000 & 0.50000 & 0.08271 & 0.28759 & 0.99253 \\
(0 - 1M)/8M &           0  & 1,000,000 & 0.50018 & 0.08331 & 0.28863 & 0.99972 \\
(0 - 8M)/8M &           0  & 8,000,000 & 0.49997 & 0.08271 & 0.28760 & 0.99257 \\
(0 - 1M)/8M & (0 - 1M)/8M  & 1,000,000 & 0.49998 & 0.08328 & 0.28858 & 0.99932 \\
(0 - 8M)/8M & (0 - 8M)/8M  & 8,000,000 & 0.50002 & 0.08271 & 0.28759 & 0.99252 \\
\hline
\end{tabular}
\caption{statistical results for \emph{maj2random} $(\mathtt{NROUNDS}=2)$.}
\end{table}

\begin{table}[h]
\centering
\small
\begin{tabular}{llrrrrr}
\hline
\textbf{x-coord} & \textbf{y-coord} & \textbf{count} & \textbf{mean} &
\textbf{variance} & \textbf{std-dev} & \textbf{$\chi^2$ ratio} \\
\hline
          0 & (0 - 1M)/8M  & 1,000,000 & 0.49993 & 0.08331 & 0.28863 & 0.99968 \\
          0 & (0 - 8M)/8M  & 8,000,000 & 0.50001 & 0.08271 & 0.28760 & 0.99258 \\
(0 - 1M)/8M &           0  & 1,000,000 & 0.49979 & 0.08332 & 0.28865 & 0.99985 \\
(0 - 8M)/8M &           0  & 8,000,000 & 0.49997 & 0.08270 & 0.28757 & 0.99237 \\
(0 - 1M)/8M & (0 - 1M)/8M  & 1,000,000 & 0.50009 & 0.08328 & 0.28859 & 0.99939 \\
(0 - 8M)/8M & (0 - 8M)/8M  & 8,000,000 & 0.50002 & 0.08273 & 0.28762 & 0.99273 \\
\hline
\end{tabular}
\caption{statistical results for \emph{maj2random} $(\mathtt{NROUNDS}=4)$.}
\end{table}

\begin{table}[h]
\centering
\small
\begin{tabular}{llrrrrr}
\hline
\textbf{x-coord} & \textbf{y-coord} & \textbf{count} & \textbf{mean} &
\textbf{variance} & \textbf{std-dev} & \textbf{$\chi^2$ ratio} \\
\hline
          0 & (0 - 1M)/8M  & 1,000,000 & 0.50000 & 0.08322 & 0.28848 & 0.99868 \\
          0 & (0 - 8M)/8M  & 8,000,000 & 0.50002 & 0.08272 & 0.28761 & 0.99261 \\
(0 - 1M)/8M &           0  & 1,000,000 & 0.49979 & 0.08329 & 0.28860 & 0.99950 \\
(0 - 8M)/8M &           0  & 8,000,000 & 0.49996 & 0.08272 & 0.28762 & 0.99268 \\
(0 - 1M)/8M & (0 - 1M)/8M  & 1,000,000 & 0.50008 & 0.08322 & 0.28847 & 0.99860 \\
(0 - 8M)/8M & (0 - 8M)/8M  & 8,000,000 & 0.49991 & 0.08270 & 0.28758 & 0.99244 \\
\hline
\end{tabular}
\caption{statistical results for \emph{maj2random} $(\mathtt{NROUNDS}=8)$.}
\end{table}

\vspace{1em}

\subsection*{conclusion}

\noindent
\emph{maj2random} exhibits uniform statistical properties and minimal
computational overhead, and its measured statistics are virtually
identical for $\mathtt{NROUNDS} \in \{2, 4, 8\}$, showing that
$\mathtt{NROUNDS}=2$ is sufficient to maintain high-quality uniform output.
given its low computational cost,  \emph{maj2random} is well-suited for
real-time graphics and procedural content generation.  \emph{note:} while
it preserves high-quality mixing, it is not intended for cryptography
applications due to the reduced number of rounds.
\\
\par\noindent
if the vector pipeline contains a \texttt{VSHA256RNDS2} instruction, then
\emph{maj2random} is an ideal way to compute high-quality deterministic
random values from texture coordinates.

\end{document}
\end{array}
